% Preview of the network-models figures on a thesis-width page (15.5 cm text). Compile twice (the shared legends use
% `legend to name`):  pdflatex preview && pdflatex preview
% In a LaTeX project: copy this folder, % Styles for the SPICE vs Spingarn figures fig_runtime.tex, fig_speedup.tex, fig_iterations.tex, fig_per_iteration.tex
% and fig_memory.tex. Needs \usepackage{pgfplots} and \usepgfplotslibrary{groupplots} (compat 1.18 tested). \input this
% file once in the preamble, then \input each figure inside a figure environment (preview.tex does both).
%
% Data: the CSVs written by `python network_models_latex.py` (repo root), one row per array size N (N x N) and one column
% per series, with _min, _max and _n (converged runs) beside each median; nan where no median exists, which pgfplots
% leaves out. runtime.csv and iterations.csv also hold _runs, _timeout and _capped (the 1600 s cap where half or more of
% the runs timed out, else nan). Series: mem_{spice,zb} (memristor alone), ideal_{spice,zb} (1D1R, ideal diode on its
% own branch), shockley_{spice,zb} (1D1R, Shockley diode on its own branch). zb is Spingarn with gamma = R_min on the
% memristors and, on 1D1R, the diode port sqrt(R_min R_max).
%
% \NMdata is the folder holding the CSVs, relative to the main .tex file; define it before \input-ing this file if the
% CSVs are elsewhere, e.g. \newcommand\NMdata{figures/network_models}.
\providecommand\NMdata{.}
\providecommand\NMcap{1600}% the runtime cap of the sweeps, in seconds

% Series colours from the repo's plotting.py (spice slot 1, spingarn slot 0), checked with the dataviz palette validator
% on white (the pair passes CVD and normal vision). Each series also has its own marker.
\definecolor{nmSpice}{HTML}{EB6834}
\definecolor{nmZeroBias}{HTML}{2A78D6}
\definecolor{nmInk}{HTML}{52514E}
\definecolor{nmGrid}{HTML}{E6E5E0}

\pgfplotsset{
  % table options: \addplot[nm spice] table[nm csv, y=mem_spice] {\NMdata/runtime.csv};
  table/nm csv/.style={col sep=comma, x=N},
  nm series/.style={color=#1, line width=1pt, mark options={fill=#1, draw=white, line width=0.4pt, solid}},
  nm spice/.style={nm series=nmSpice, mark=*, mark size=2pt},
  nm zero bias/.style={nm series=nmZeroBias, mark=square*, mark size=1.8pt},
  % min-max over the seeds as error bars:
  %   \addplot[nm spice, nm bars] table[nm csv, y=mem_spice, nm range=mem_spice] {\NMdata/runtime.csv};
  nm bars/.style={error bars/.cd, y dir=both, y explicit, error bar style={line width=0.6pt, solid},
    error mark options={rotate=90, mark size=1.6pt, line width=0.6pt}},
  table/nm range/.style={y error plus expr=\thisrow{#1_max}-\thisrow{#1},
    y error minus expr=\thisrow{#1}-\thisrow{#1_min}},
  % a run over the cap: a cross at the cap, labelled; \addplot[nm timeout=nmSpice] table[nm csv, y=X_capped] {...};
  nm timeout/.style={only marks, mark=x, mark size=3.2pt, line width=1.3pt, color=#1, forget plot,
    nodes near coords={timed out}, every node near coord/.append style={font=\scriptsize, text=#1, anchor=south east,
    xshift=-1pt}},
  nm cap/.style={draw=nmInk!70, line width=0.6pt, dash pattern=on 3pt off 2pt, forget plot},
  % a lower bound on the speedup (SPICE over the cap): an open triangle in Spingarn's colour
  nm bound/.style={only marks, mark=triangle, mark size=2.8pt, color=nmZeroBias,
    mark options={draw=nmZeroBias, fill=white, line width=0.9pt, solid}},
  nm reference/.style={draw=nmInk!55, line width=0.6pt, dashed, forget plot},
  % one panel; width/height are per panel (three panels plus labels fill ~15.5 cm)
  nm axis/.style={
    width=5.3cm, height=5cm, scale only axis=false,
    xmode=log, ymode=log,
    xtick={16,64,256,1024}, xticklabels={16,64,256,1024}, minor xtick={32,128,512},
    xmin=12, xmax=1400,
    yminorticks=false,
    xlabel={$N$ ($N\times N$ array)},
    grid=major, major grid style={draw=nmGrid, line width=0.4pt},
    axis lines=left, axis line style={draw=nmInk!60, -},
    tick style={draw=nmInk!60}, tick label style={font=\footnotesize, text=nmInk},
    label style={font=\small, text=nmInk}, title style={font=\small, text=black},
    unbounded coords=jump,
  },
  % the three-panel layout; groupplot reads group style only from its own options, so use it as
  % \begin{groupplot}[group style={nm group style}, nm axis, ...]
  group/nm group style/.style={group size=3 by 1, horizontal sep=0.45cm, y descriptions at=edge left},
  nm legend/.style={legend columns=3, legend style={draw=none, font=\small, text=nmInk,
    /tikz/every even column/.append style={column sep=0.8em}}},
}
 in the preamble (after pgfplots and groupplots;
% define \NMdata first if the CSVs are not next to the main file), and \input{fig_...} inside a figure environment.
\documentclass[11pt]{article}
\usepackage[a4paper, textwidth=15.5cm, top=2cm, bottom=2cm]{geometry}
\usepackage{pgfplots}
\usepgfplotslibrary{groupplots}
\pgfplotsset{compat=1.18}
% Styles for the SPICE vs Spingarn figures fig_runtime.tex, fig_speedup.tex, fig_iterations.tex, fig_per_iteration.tex
% and fig_memory.tex. Needs \usepackage{pgfplots} and \usepgfplotslibrary{groupplots} (compat 1.18 tested). \input this
% file once in the preamble, then \input each figure inside a figure environment (preview.tex does both).
%
% Data: the CSVs written by `python network_models_latex.py` (repo root), one row per array size N (N x N) and one column
% per series, with _min, _max and _n (converged runs) beside each median; nan where no median exists, which pgfplots
% leaves out. runtime.csv and iterations.csv also hold _runs, _timeout and _capped (the 1600 s cap where half or more of
% the runs timed out, else nan). Series: mem_{spice,zb} (memristor alone), ideal_{spice,zb} (1D1R, ideal diode on its
% own branch), shockley_{spice,zb} (1D1R, Shockley diode on its own branch). zb is Spingarn with gamma = R_min on the
% memristors and, on 1D1R, the diode port sqrt(R_min R_max).
%
% \NMdata is the folder holding the CSVs, relative to the main .tex file; define it before \input-ing this file if the
% CSVs are elsewhere, e.g. \newcommand\NMdata{figures/network_models}.
\providecommand\NMdata{.}
\providecommand\NMcap{1600}% the runtime cap of the sweeps, in seconds

% Series colours from the repo's plotting.py (spice slot 1, spingarn slot 0), checked with the dataviz palette validator
% on white (the pair passes CVD and normal vision). Each series also has its own marker.
\definecolor{nmSpice}{HTML}{EB6834}
\definecolor{nmZeroBias}{HTML}{2A78D6}
\definecolor{nmInk}{HTML}{52514E}
\definecolor{nmGrid}{HTML}{E6E5E0}

\pgfplotsset{
  % table options: \addplot[nm spice] table[nm csv, y=mem_spice] {\NMdata/runtime.csv};
  table/nm csv/.style={col sep=comma, x=N},
  nm series/.style={color=#1, line width=1pt, mark options={fill=#1, draw=white, line width=0.4pt, solid}},
  nm spice/.style={nm series=nmSpice, mark=*, mark size=2pt},
  nm zero bias/.style={nm series=nmZeroBias, mark=square*, mark size=1.8pt},
  % min-max over the seeds as error bars:
  %   \addplot[nm spice, nm bars] table[nm csv, y=mem_spice, nm range=mem_spice] {\NMdata/runtime.csv};
  nm bars/.style={error bars/.cd, y dir=both, y explicit, error bar style={line width=0.6pt, solid},
    error mark options={rotate=90, mark size=1.6pt, line width=0.6pt}},
  table/nm range/.style={y error plus expr=\thisrow{#1_max}-\thisrow{#1},
    y error minus expr=\thisrow{#1}-\thisrow{#1_min}},
  % a run over the cap: a cross at the cap, labelled; \addplot[nm timeout=nmSpice] table[nm csv, y=X_capped] {...};
  nm timeout/.style={only marks, mark=x, mark size=3.2pt, line width=1.3pt, color=#1, forget plot,
    nodes near coords={timed out}, every node near coord/.append style={font=\scriptsize, text=#1, anchor=south east,
    xshift=-1pt}},
  nm cap/.style={draw=nmInk!70, line width=0.6pt, dash pattern=on 3pt off 2pt, forget plot},
  % a lower bound on the speedup (SPICE over the cap): an open triangle in Spingarn's colour
  nm bound/.style={only marks, mark=triangle, mark size=2.8pt, color=nmZeroBias,
    mark options={draw=nmZeroBias, fill=white, line width=0.9pt, solid}},
  nm reference/.style={draw=nmInk!55, line width=0.6pt, dashed, forget plot},
  % one panel; width/height are per panel (three panels plus labels fill ~15.5 cm)
  nm axis/.style={
    width=5.3cm, height=5cm, scale only axis=false,
    xmode=log, ymode=log,
    xtick={16,64,256,1024}, xticklabels={16,64,256,1024}, minor xtick={32,128,512},
    xmin=12, xmax=1400,
    yminorticks=false,
    xlabel={$N$ ($N\times N$ array)},
    grid=major, major grid style={draw=nmGrid, line width=0.4pt},
    axis lines=left, axis line style={draw=nmInk!60, -},
    tick style={draw=nmInk!60}, tick label style={font=\footnotesize, text=nmInk},
    label style={font=\small, text=nmInk}, title style={font=\small, text=black},
    unbounded coords=jump,
  },
  % the three-panel layout; groupplot reads group style only from its own options, so use it as
  % \begin{groupplot}[group style={nm group style}, nm axis, ...]
  group/nm group style/.style={group size=3 by 1, horizontal sep=0.45cm, y descriptions at=edge left},
  nm legend/.style={legend columns=3, legend style={draw=none, font=\small, text=nmInk,
    /tikz/every even column/.append style={column sep=0.8em}}},
}


\begin{document}

\begin{figure}[p]
  \centering
  % Runtime vs array size, SPICE vs Spingarn on the three crossbar models: medians over 10 seeds with min-max bars, the PC
% (DESKTOP-EFQTE92). A timed-out run counts as above the cap in the median; where half or more timed out the median is
% missing and a cross marks the cap. Needs network_models_style.tex; data runtime.csv (s).
\begin{tikzpicture}
\begin{groupplot}[group style={nm group style}, nm axis, ylabel={runtime (s)},
  ymin=1e-3, ymax=1e4, ytick={1e-3,1e-2,1e-1,1,1e1,1e2,1e3,1e4}]
\nextgroupplot[title={Memristor}, legend to name=nmRuntimeLegend, nm legend]
  \addplot[nm spice, nm bars]     table[nm csv, y=mem_spice, nm range=mem_spice] {\NMdata/runtime.csv};
    \addlegendentry{SPICE (ngspice, KLU)}
  \addplot[nm zero bias, nm bars] table[nm csv, y=mem_zb, nm range=mem_zb]       {\NMdata/runtime.csv};
    \addlegendentry{Spingarn, zero-bias $\gamma$}
  \addlegendimage{nm cap}\addlegendentry{\NMcap\,s runtime cap}
  \addplot[nm cap, domain=12:1400, samples=2] {\NMcap};
  \addplot[nm timeout=nmSpice]    table[nm csv, y=mem_spice_capped] {\NMdata/runtime.csv};
  \addplot[nm timeout=nmZeroBias] table[nm csv, y=mem_zb_capped]    {\NMdata/runtime.csv};
\nextgroupplot[title={1D1R, ideal diode}]
  \addplot[nm spice, nm bars]     table[nm csv, y=ideal_spice, nm range=ideal_spice] {\NMdata/runtime.csv};
  \addplot[nm zero bias, nm bars] table[nm csv, y=ideal_zb, nm range=ideal_zb]       {\NMdata/runtime.csv};
  \addplot[nm cap, domain=12:1400, samples=2] {\NMcap};
  \addplot[nm timeout=nmSpice]    table[nm csv, y=ideal_spice_capped] {\NMdata/runtime.csv};
  \addplot[nm timeout=nmZeroBias] table[nm csv, y=ideal_zb_capped]    {\NMdata/runtime.csv};
\nextgroupplot[title={1D1R, Shockley diode}]
  \addplot[nm spice, nm bars]     table[nm csv, y=shockley_spice, nm range=shockley_spice] {\NMdata/runtime.csv};
  \addplot[nm zero bias, nm bars] table[nm csv, y=shockley_zb, nm range=shockley_zb]       {\NMdata/runtime.csv};
  \addplot[nm cap, domain=12:1400, samples=2] {\NMcap};
  \addplot[nm timeout=nmSpice]    table[nm csv, y=shockley_spice_capped] {\NMdata/runtime.csv};
  \addplot[nm timeout=nmZeroBias] table[nm csv, y=shockley_zb_capped]    {\NMdata/runtime.csv};
\end{groupplot}
\node[anchor=south, yshift=2pt] at (group c2r1.outer north) {\pgfplotslegendfromname{nmRuntimeLegend}};
\end{tikzpicture}

  \caption{Runtime against array size for SPICE (ngspice-47, KLU) and Spingarn with fixed port resistances
  ($\alpha = 1$ on the memristor, $0.8$ on 1D1R; CHOLMOD with AMD ordering), median over three random circuits.
  Zero-bias $\gamma$: $\gamma = R_{\min}$ per memristor; sector $\gamma$: $\sqrt{R_{\min}R_{\max}}$. On 1D1R the diode
  is on its own branch with port resistance $\sqrt{R_{\min}R_{\max}}$. SPICE writes the ideal diode as
  $\mathrm{uramp}(V)/1\,\mathrm{m\Omega}$ and the Shockley diode ($I_s = 10^{-9}$\,A, $n = 1$) as a behavioural
  $\exp$ source; on that case every SPICE run needed gmin stepping, and its $512\times512$ point is a single run.
  SPICE was not run at $1024\times1024$, where ngspice's KLU fails.}
\end{figure}

\begin{figure}[p]
  \centering
  % SPICE's runtime over Spingarn's (> 1: Spingarn faster), from the medians in runtime.csv; only where SPICE has data.
% Needs network_models_style.tex.
\begin{tikzpicture}
\begin{groupplot}[group style={nm group style}, nm axis, ylabel={SPICE runtime / Spingarn runtime},
  xmax=700, ymin=0.2, ymax=60, ytick={0.25,0.5,1,2,4,8,16,32}, yticklabels={0.25,0.5,1,2,4,8,16,32}]
\nextgroupplot[title={Memristor}, legend to name=nmSpeedupLegend, nm legend]
  \addplot[nm reference, domain=12:700, samples=2] {1};
  \addplot[nm zero bias] table[nm csv, y expr=\thisrow{mem_spice}/\thisrow{mem_zb}]     {\NMdata/runtime.csv};
    \addlegendentry{Spingarn, zero-bias $\gamma$}
  \addplot[nm sector]    table[nm csv, y expr=\thisrow{mem_spice}/\thisrow{mem_sector}] {\NMdata/runtime.csv};
    \addlegendentry{Spingarn, sector $\gamma$}
  \node[font=\scriptsize, text=nmInk, anchor=north west] at (rel axis cs:0.03,0.97) {Spingarn faster $\uparrow$};
  \node[font=\scriptsize, text=nmInk, anchor=south west] at (rel axis cs:0.03,0.03) {SPICE faster $\downarrow$};
\nextgroupplot[title={1D1R, ideal diode}]
  \addplot[nm reference, domain=12:700, samples=2] {1};
  \addplot[nm zero bias] table[nm csv, y expr=\thisrow{ideal_spice}/\thisrow{ideal_zb}] {\NMdata/runtime.csv};
\nextgroupplot[title={1D1R, Shockley diode}]
  \addplot[nm reference, domain=12:700, samples=2] {1};
  \addplot[nm zero bias] table[nm csv, y expr=\thisrow{shockley_spice}/\thisrow{shockley_zb}] {\NMdata/runtime.csv};
\end{groupplot}
\node[anchor=south, yshift=2pt] at (group c2r1.outer north) {\pgfplotslegendfromname{nmSpeedupLegend}};
\end{tikzpicture}

  \caption{SPICE's runtime divided by Spingarn's, from the medians of the runtime figure.}
\end{figure}

\begin{figure}[p]
  \centering
  % Iterations vs array size: Spingarn's iterations and SPICE's Newton iterations (fallback stages such as gmin stepping
% included); medians over 3 seeds. Needs network_models_style.tex; data iterations.csv.
\begin{tikzpicture}
\begin{groupplot}[group style={nm group style}, nm axis, ylabel={iterations},
  ymin=2, ymax=1e4, ytick={1,1e1,1e2,1e3,1e4}]
\nextgroupplot[title={Memristor}, legend to name=nmIterationsLegend, nm legend]
  \addplot[nm spice]     table[nm csv, y=mem_spice]  {\NMdata/iterations.csv}; \addlegendentry{SPICE (Newton)}
  \addplot[nm zero bias] table[nm csv, y=mem_zb]     {\NMdata/iterations.csv}; \addlegendentry{Spingarn, zero-bias $\gamma$}
  \addplot[nm sector]    table[nm csv, y=mem_sector] {\NMdata/iterations.csv}; \addlegendentry{Spingarn, sector $\gamma$}
\nextgroupplot[title={1D1R, ideal diode}]
  \addplot[nm spice]     table[nm csv, y=ideal_spice] {\NMdata/iterations.csv};
  \addplot[nm zero bias] table[nm csv, y=ideal_zb]    {\NMdata/iterations.csv};
\nextgroupplot[title={1D1R, Shockley diode}]
  \addplot[nm spice]     table[nm csv, y=shockley_spice] {\NMdata/iterations.csv};
  \addplot[nm zero bias] table[nm csv, y=shockley_zb]    {\NMdata/iterations.csv};
\end{groupplot}
\node[anchor=south, yshift=2pt] at (group c2r1.outer north) {\pgfplotslegendfromname{nmIterationsLegend}};
\end{tikzpicture}

  \caption{Iterations: Spingarn's, and SPICE's Newton iterations including fallback stages (gmin stepping).}
\end{figure}

\begin{figure}[p]
  \centering
  % Time per iteration vs array size: SPICE per Newton iteration (ngspice's analysis time less KLU's reorder, over its
% iterations) and Spingarn per iteration (the iteration loop over its iterations); one-time costs excluded (one_time.csv).
% Medians over 3 seeds. Needs network_models_style.tex; data per_iteration.csv (ms).
\begin{tikzpicture}
\begin{groupplot}[group style={nm group style}, nm axis, ylabel={time per iteration (ms)},
  ymin=0.05, ymax=5e4, ytick={1e-1,1,1e1,1e2,1e3,1e4}]
\nextgroupplot[title={Memristor}, legend to name=nmPerIterationLegend, nm legend]
  \addplot[nm spice]     table[nm csv, y=mem_spice]  {\NMdata/per_iteration.csv}; \addlegendentry{SPICE (per Newton iteration)}
  \addplot[nm zero bias] table[nm csv, y=mem_zb]     {\NMdata/per_iteration.csv}; \addlegendentry{Spingarn, zero-bias $\gamma$}
  \addplot[nm sector]    table[nm csv, y=mem_sector] {\NMdata/per_iteration.csv}; \addlegendentry{Spingarn, sector $\gamma$}
\nextgroupplot[title={1D1R, ideal diode}]
  \addplot[nm spice]     table[nm csv, y=ideal_spice] {\NMdata/per_iteration.csv};
  \addplot[nm zero bias] table[nm csv, y=ideal_zb]    {\NMdata/per_iteration.csv};
\nextgroupplot[title={1D1R, Shockley diode}]
  \addplot[nm spice]     table[nm csv, y=shockley_spice] {\NMdata/per_iteration.csv};
  \addplot[nm zero bias] table[nm csv, y=shockley_zb]    {\NMdata/per_iteration.csv};
\end{groupplot}
\node[anchor=south, yshift=2pt] at (group c2r1.outer north) {\pgfplotslegendfromname{nmPerIterationLegend}};
\end{tikzpicture}

  \caption{Time per iteration, one-time costs (netlist load and KLU reordering; Spingarn's assembly and
  factorization) excluded. On the memristor the two Spingarn variants coincide: they factorize the same matrix
  and do the same work per iteration.}
\end{figure}

\begin{figure}[p]
  \centering
  % Peak memory vs array size: the peak working set of the process each run executed in (Python, numpy and, for SPICE,
% ngspice included; ~0.1 GB at the smallest sizes is that baseline). Medians over converged runs of 10 seeds. Needs
% network_models_style.tex; data memory.csv (GB).
\begin{tikzpicture}
\begin{groupplot}[group style={nm group style}, nm axis, ylabel={peak memory (GB)},
  ymin=0.05, ymax=20, ytick={0.1,1,10}, yticklabels={0.1,1,10}]
\nextgroupplot[title={Memristor}, legend to name=nmMemoryLegend, nm legend]
  \addplot[nm spice]     table[nm csv, y=mem_spice] {\NMdata/memory.csv}; \addlegendentry{SPICE (ngspice, KLU)}
  \addplot[nm zero bias] table[nm csv, y=mem_zb]    {\NMdata/memory.csv}; \addlegendentry{Spingarn, zero-bias $\gamma$}
\nextgroupplot[title={1D1R, ideal diode}]
  \addplot[nm spice]     table[nm csv, y=ideal_spice] {\NMdata/memory.csv};
  \addplot[nm zero bias] table[nm csv, y=ideal_zb]    {\NMdata/memory.csv};
\nextgroupplot[title={1D1R, Shockley diode}]
  \addplot[nm spice]     table[nm csv, y=shockley_spice] {\NMdata/memory.csv};
  \addplot[nm zero bias] table[nm csv, y=shockley_zb]    {\NMdata/memory.csv};
\end{groupplot}
\node[anchor=south, yshift=2pt] at (group c2r1.outer north) {\pgfplotslegendfromname{nmMemoryLegend}};
\end{tikzpicture}

  \caption{Peak memory of the process each run executed in. On the memristor the two Spingarn variants
  coincide.}
\end{figure}

\end{document}
